% Bench layout: one shield for the read, one pin for the toggle, one marker.
\begin{tikzpicture}[font=\sffamily\scriptsize]
\node[board, minimum width=52mm, minimum height=28mm] (nuc) at (0,0) {};
\node[text=white, font=\sffamily\bfseries\small] at (-0.55,-1.05) {NUCLEO-H7A3ZI-Q};
\node[pinrow, minimum width=42mm] at (0,1.25) {};
\node[pinrow, minimum width=42mm] at (0,-0.55) {};

\node[hat, minimum width=38mm, minimum height=15mm] (sh) at (0.1,0.5) {};
\node[text=white, font=\sffamily\bfseries\small] at (0.1,0.85) {one sensor shield};
\node[text=white, font=\tiny, align=center] at (0.1,0.32)
  {either shield serves: the operation\\under test is a bus transaction};

\node[font=\tiny, align=right, anchor=south east] at (-3.05,0.95) {ST-LINK\\micro USB};
\node[draw=black!70, fill=gray!60, minimum width=4mm, minimum height=6mm] (usb) at (-2.75,0.6) {};
\node[ext, left=16mm of nuc, text width=22mm, minimum height=14mm] (pc)
  {Host PC\\{\tiny two toolchains\\one comparison script}};
\draw[usbcable] (usb.west) to[out=180,in=0] (pc.east);

\node[ext, right=18mm of nuc, text width=25mm, minimum height=17mm] (ppk)
  {nRF-PPK2\\{\tiny the marker recorder, and charge per operation if wanted}};
\draw[cable, draw=orange!70!black] (2.6,0.2) -- node[above, font=\tiny]{marker} (ppk.west);
\draw[cable, draw=blue!50!black] (2.6,-0.25)
  -- node[below, font=\tiny]{toggled pin} ($(ppk.west)+(0,-0.45)$);

\node[umlnote, text width=88mm, anchor=north west] at (-4.6,-2.0)
  {The toggled pin drives nothing but a logic input, so there is no load to
   argue about. There is no oscilloscope on this bench and a toggle at these
   rates cannot be observed externally with what is here, so the cycle counter
   is the instrument, its own cost is measured and subtracted in the harness,
   and the substitution is stated rather than implied. 3.3 V logic only.};
\end{tikzpicture}
